% Phụ lục A

\chapter{Định nghĩa các nhãn cảm xúc} % Tên của phụ lục

\label{AppendixA} % Để trích dẫn chương này ở chỗ nào đó trong bài, hãy sử dụng lệnh \ref{AppendixA} 

%----------------------------------------------------------------------------------------

Hệ thống nhãn cảm xúc trong bộ dữ liệu GoEmotions bao gồm 27 danh mục sắc thái cảm xúc chi tiết và 1 nhãn trung tính, được định nghĩa cụ thể như sau:

\begin{enumerate}
	\item \textbf{admiration} \twemoji{1f44f}: NGƯỠNG MỘ (Thể hiện việc tìm thấy một điều gì đó ấn tượng hoặc đáng trân trọng, kính phục).
	
	\item \textbf{amusement} \twemoji{1f602}: THÍCH THÚ (Thể hiện việc tìm thấy một điều gì đó hài hước, vui nhộn hoặc mang tính giải trí).
	
	\item \textbf{anger} \twemoji{1f621}: GIẬN DỮ (Cảm giác không hài lòng mạnh mẽ, bất bình hoặc mang tính phản kháng, đối địch).
	
	\item \textbf{annoyance} \twemoji{1f612}: KHÓ CHỊU (Trạng thái tức giận nhẹ, bực mình hoặc kích động không thoải mái).
	
	\item \textbf{approval} \twemoji{1f44d}: TÁN THÀNH (Việc sở hữu hoặc biểu đạt một quan điểm, ý kiến thuận chiều, ủng hộ).
	
	\item \textbf{caring} \twemoji{1f970}: QUAN TÂM (Biểu lộ sự tử tế, lòng trắc ẩn và sự lo lắng, chú ý dành cho người khác).
	
	\item \textbf{confusion} \twemoji{1f928}: BỐI RỐI (Sự thiếu thấu hiểu, mơ hồ hoặc trạng thái không chắc chắn về một vấn đề).
	
	\item \textbf{curiosity} \twemoji{1f914}: TÒ MÒ (Mong muốn mạnh mẽ được biết, được khám phá hoặc học hỏi một điều gì đó).
	
	\item \textbf{desire} \twemoji{1f97a}: KHAO KHÁT (Cảm giác mạnh mẽ muốn có được một thứ gì đó hoặc mong ước một sự việc cụ thể xảy ra).
	
	\item \textbf{disappointment} \twemoji{1f61e}: THẤT VỌNG (Sắc thái buồn bã hoặc không hài lòng gây ra bởi việc không đạt được kỳ vọng hay hy vọng).
	
	\item \textbf{disapproval} \twemoji{1f44e}: PHẢN ĐỐI (Việc sở hữu hoặc biểu đạt một quan điểm, ý kiến không thuận chiều, bất đồng).
	
	\item \textbf{disgust} \twemoji{1f92e}: GHÊ TỞM (Sự phẫn nộ hoặc phản đối mạnh mẽ được khơi dậy bởi một điều gì đó khó chịu hoặc phản cảm).
	
	\item \textbf{embarrassment} \twemoji{1f633}: NGƯỢNG NGÙNG (Trạng thái tự ý thức tiêu cực, xấu hổ, ngượng nghịu hoặc lúng túng trong ngữ cảnh xã hội).
	
	\item \textbf{excitement} \twemoji{1f929}: HÀO HỨNG (Cảm giác nhiệt huyết lớn, mong đợi mãnh liệt và tràn đầy năng lượng tích cực).
	
	\item \textbf{fear} \twemoji{1f628}: SỢ HÃI (Trạng thái hoảng sợ, lo lắng sâu sắc hoặc e ngại trước một mối nguy hại).
	
	\item \textbf{gratitude} \twemoji{1f64f}: BIẾT ƠN (Cảm giác biết ơn, trân trọng và đánh giá cao sự giúp đỡ hoặc một điều tốt đẹp nhận được).
	
	\item \textbf{grief} \twemoji{1f62d}: ĐAU BUỒN (Nỗi đau khổ dữ dội, tột cùng, đặc biệt là nỗi đau gây ra bởi sự ra đi của một ai đó).
	
	\item \textbf{joy} \twemoji{1f600}: VUI SƯỚNG (Cảm giác cực kỳ hạnh phúc, thỏa mãn và tràn ngập niềm vui).
	
	\item \textbf{love} \twemoji{2764}: YÊU THƯƠNG (Cảm xúc tích cực mạnh mẽ thể hiện sự gắn kết, quý mến và tình cảm sâu sắc).
	
	\item \textbf{nervousness} \twemoji{1f62c}: LO LẮNG (Trạng thái bồn chồn, e sợ, căng thẳng hoặc bất an trước một sự kiện).
	
	\item \textbf{optimism} \twemoji{1f91e}: LẠC QUAN (Sự hy vọng và niềm tin tưởng mạnh mẽ vào tương lai hoặc sự thành công của một công việc).
	
	\item \textbf{pride} \twemoji{1f60e}: TỰ HÀO (Niềm vui lớn hoặc sự thỏa mãn sâu sắc đến từ thành tựu của chính mình hoặc của những người gắn bó).
	
	\item \textbf{realization} \twemoji{1f4a1}: NHẬN RA (Trạng thái bắt đầu nhận thức, thấu hiểu hoặc đột nhiên thức tỉnh về một sự thật nào đó).
	
	\item \textbf{relief} \twemoji{1f60c}: NHẸ NHÕM (Sự trấn tĩnh, an tâm và thư giãn sau khi được giải phóng khỏi trạng thái lo âu hoặc đau khổ).
	
	\item \textbf{remorse} \twemoji{1f625}: HỐI HẬN (Cảm giác hối tiếc sâu sắc hoặc dằn vặt, có lỗi về một hành vi sai lầm đã qua).
	
	\item \textbf{sadness} \twemoji{1f622}: BUỒN BÃ (Nỗi đau đớn về mặt cảm xúc, tâm trạng sầu não, chán nản).
	
	\item \textbf{surprise} \twemoji{1f62e}: NGẠC NHIÊN (Cảm giác kinh ngạc, sửng sốt hoặc giật mình gây ra bởi một điều gì đó hoàn toàn bất ngờ).
	
	\item \textbf{neutral} \twemoji{1f610}: TRUNG TÍNH (Văn bản thuần túy thông tin thực tế, không chứa bất kỳ sắc thái biểu cảm hay cảm xúc nào).
\end{enumerate}